% Labb 3, parts 1.1, 1.2 and 2, from labs/lab3/a_first_steps.md. Each part is a \labtask, and its
% Mål, Bakgrund, Programmet, Uppgifter and Frågor are \task headings. The program given in part
% 1.1 and the four faulty lines in part 1.2 are the handout's, and are printed as the guide
% prints them.

\section{Del 1.1, 1.2 och 2: Första stegen}
\label{lab3:sec:a}
Delarna hör till kapitel~\ref{c7:ch}. Läs \secref{c7:app:b} och \secref{c7:app:c} innan ni börjar,
och ha bilaga~\ref{app:studio} om Microchip Studio öppen bredvid.

\labtask{Del 1.1}{Microchip Studio: skapa, skriva, spara, assemblera, simulera}%
\phantomsection\label{lab3:d1-1}
\task{Mål}
Skapa ett assemblerprojekt, skriva ett program, bygga det och stega igenom det i simulatorn.

\task{Bakgrund}
Arbetsgången, skriv, bygg, förutsäg, stega och jämför, beskrivs i \secref{c7:c1}. Instruktionerna
i programmet står i \secref{c7:b5}.

\task{Programmet}
\begin{asmcode}
.include "m328Pdef.inc"

.org 0x0000
    rjmp main

main:
    ldi r16, 10                     ; r16 = 10
    ldi r17, 20                     ; r17 = 20
    mov r18, r16                    ; r18 = r16
    inc r18                         ; r18 = r18 + 1
    inc r18                         ; r18 = r18 + 1
    dec r17                         ; r17 = r17 - 1

end:
    rjmp end
\end{asmcode}

\task{Uppgifter}
\begin{enumerate}
  \item Skapa ett projekt med namnet \code{lab3_01_1} enligt \secref{studio:2}.
  \item Ersätt innehållet i \code{main.asm} med programmet ovan. \textbf{Skriv av det} i stället
    för att klistra in: det är så ni lär er hur en rad är uppbyggd. Kommentarerna kan ni hoppa
    över.
  \item Spara med \textbf{Ctrl+S}.
  \item Bygg med \textbf{F7}. Läs utskriften i fönstret Output: står det \code{0 errors}? Om inte,
    rätta raden Error List pekar på och bygg om.
  \item Välj simulatorn som verktyg (\secref{studio:4-1}) och starta med \textbf{Debug
    $\rightarrow$ Start Debugging and Break} (\textbf{Alt+F5}). Öppna fönstret \textbf{Processor
    Status}.
  \item \textbf{Innan ni stegar:} fyll i tabellen nedan med registrens värden efter varje
    instruktion, decimalt och hexadecimalt.

    \begin{filltable}{@{}p{8.4em}YYY@{}}
      \thead{Efter} & \thead{\code{r16}} & \thead{\code{r17}} & \thead{\code{r18}} \\
      \midrule
      \code{ldi r16, 10} & & & \fillrule
      \code{ldi r17, 20} & & & \fillrule
      \code{mov r18, r16} & & & \fillrule
      \code{inc r18} & & & \fillrule
      \code{inc r18} & & & \fillrule
      \code{dec r17} & & & \\
    \end{filltable}
  \item Stega med \textbf{F11}, en instruktion i taget, och jämför Processor Status med tabellen
    efter varje steg. Lägg märke till att registret som just ändrades visas i rött.
  \item Läs av \textbf{Program Counter} vid varje steg. Med hur mycket ökar den för varje
    instruktion?
  \item När den gula pilen står på \code{rjmp end}: vad visar \textbf{Cycle Counter}? Räkna ut
    värdet för hand först, med hjälp av \secref{c7:c5}.
\end{enumerate}

\task{Frågor}
\begin{itemize}
  \item Den gula pilen står på \code{inc r18}. Har \code{inc r18} körts ännu?
  \item Ändras \code{r16} av instruktionen \code{mov r18, r16}?
\end{itemize}

\redovisa{tabellen och simuleringen för läraren.}

\labtask{Del 1.2}{Öppna, ändra, rätta, assemblera, simulera}\phantomsection\label{lab3:d1-2}
\task{Mål}
Öppna ett befintligt projekt, ändra programmet, och rätta fel med hjälp av felmeddelandena.

\task{Bakgrund}
De vanligaste felmeddelandena, och hur de rättas, står i \secref{c7:c4}.

\task{Uppgifter}
\begin{enumerate}
  \item Stäng Microchip Studio och starta det igen. Öppna projektet från del 1.1, antingen från
    listan över senaste projekt på startsidan eller med \textbf{File $\rightarrow$ Open
    $\rightarrow$ Project/Solution...}
  \item Ändra programmet:
    \begin{itemize}
      \item \code{r16} ska laddas med 25 i stället för 10;
      \item lägg till raden \code{mov r19, r17} efter \code{dec r17}.
    \end{itemize}
  \item \textbf{Förutsäg} värdena i \code{r16}--\code{r19} när programmet når \code{end}. Bygg,
    simulera och kontrollera.
  \item Lägg nu till de här fyra raderna direkt efter \code{mov r19, r17}. De innehåller fyra fel
    med flit:
\begin{asmcode}
    ldi r16 7
    ldi r5, 3
    inc r19, 2
    rjmp ennd
\end{asmcode}
    Meningen med raderna är: ladda 7 i \code{r16}, ladda 3 i något register, öka \code{r19} med
    två, och hoppa till \code{end}.
  \item Bygg. Läs det första felet i \textbf{Error List}, dubbelklicka på det, rätta raden och bygg
    om. Fortsätt tills programmet byggs utan fel. Skriv för varje fel ned vad meddelandet sa och
    hur ni rättade det.
  \item \textbf{Förutsäg} värdena i \code{r16}--\code{r20} när programmet når \code{end}. Simulera
    och kontrollera.
\end{enumerate}

\task{Frågor}
\begin{itemize}
  \item Hur många gånger behövde ni bygga om innan alla fel var borta? Fick ni alla fyra felen på
    en gång?
  \item Varför kan \code{inc} inte öka med två på en gång?
\end{itemize}

\redovisa{det rättade programmet och de fyra felmeddelandena.}

\labtask{Del 2}{Decimalt, binärt eller hexadecimalt}\phantomsection\label{lab3:d2}
\task{Mål}
Skriva tal i ett program decimalt, binärt och hexadecimalt, och läsa dem i simulatorn.

\task{Bakgrund}
Talformaten står i \secref{c7:b4} och omvandlingen mellan talsystemen i kapitel~\ref{c1:ch}.

\task{Uppgifter}
\begin{enumerate}
  \item Skriv ett program som laddar talet 200 i \code{r16} decimalt, i \code{r17} binärt och i
    \code{r18} hexadecimalt. Räkna ut det binära och det hexadecimala skrivsättet för hand.
    \textbf{Förutsäg} hur Processor Status visar de tre registren. Simulera och kontrollera.
  \item Fyll i de tomma rutorna i tabellen \textbf{för hand}:

    \begin{filltable}{@{}p{5em}YYY@{}}
      \thead{Register} & \thead{Decimalt} & \thead{Binärt} & \thead{Hexadecimalt} \\
      \midrule
      \code{r19} & 9 & & \fillrule
      \code{r20} & & \code{0b00001111} & \fillrule
      \code{r21} & & & \code{0x10} \fillrule
      \code{r22} & 100 & & \fillrule
      \code{r23} & & \code{0b01111111} & \fillrule
      \code{r24} & & & \code{0x80} \fillrule
      \code{r25} & 255 & & \\
    \end{filltable}
  \item Lägg till en \code{ldi} per rad i programmet, som laddar registret med talet \textbf{i det
    skrivsätt tabellen ger det}: \code{ldi r19, 9}, \code{ldi r20, 0b00001111}, och så vidare.
    Simulera och kontrollera att Processor Status visar det hexadecimala värde ni räknade ut på
    varje rad.
  \item Lägg till \code{ldi r26, '0'} och \code{ldi r27, 'A'}. Vilka värden får registren? Ett
    tecken lagras som ett tal, dess ASCII-kod; det återkommer i del~18.
  \item Välj för varje rad i tabellen vilket skrivsätt som passar bäst om talet är
    \begin{itemize}
      \item ett antal varv i en loop;
      \item ett mönster av lysdioder som ska lysa.
    \end{itemize}
\end{enumerate}

\task{Frågor}
\begin{itemize}
  \item Blir maskinkoden olika för \code{ldi r16, 200} och \code{ldi r16, 0xC8}? Titta i listfilen
    om ni är osäkra (\secref{c7:b7}).
\end{itemize}

\redovisa{tabellen och simuleringen.}
